% GENERATED FILE. Edit canonical lessons, then run npm run generate:books.
% canonical-source-hash: fnv1a64:4e27078ecad569f0
% canonical-lessons: PA-C200-vegetable-return, PA-C200-sabzi, PA-C200-sabzi-review, PA-C200-vegetable-label, PA-C200-payoff

\chapter{This Is a Vegetable}
\label{ch:pa-a2-sabzi}
\hlchaptermodality{200}

\begin{chapteropening}
\textbf{By the end of this chapter:} \emph{I can recognise and use the noun for a vegetable or cooked vegetable dish in one bounded identity frame.}

Recognise and produce the vegetable noun in one bounded identity frame through listening, speaking, reading, and writing.
\end{chapteropening}

\section[\texorpdfstring{\pa{ਇਹ} · \pa{ਹੈ} (eh; hai)}{ਇਹ · ਹੈ (eh; hai)}]{Return to this and is}

\label{lesson:PA-C200-vegetable-return}



\begin{quote}
\emph{Hear:} \emph{eh; hai}; YOU SAY: \textbf{\pa{ਇਹ}; \pa{ਹੈ}}
\end{quote}

\begin{grammarlens}[title={Grammar Lens}]
Keep the familiar pointer first and the familiar verb last. The middle slot stays empty until the next lesson owns its new word.
\end{grammarlens}

\subsection*{Your turn}
\begin{itemize}
\raggedright
  \item \emph{Hear:} \emph{eh}; YOU SAY: \textbf{\pa{ਇਹ}}
  \item \emph{Hear:} \emph{hai}; YOU SAY: \textbf{\pa{ਹੈ}}
  \item Say the first word, pause for the empty middle, then say the last.
\end{itemize}

\begin{tcolorbox}[breakable,colback=teal!4,colframe=teal!35!black,title={Before you move on}]
Produce the two familiar words around the empty middle slot. This opener owns no new word.
\end{tcolorbox}
\section[\texorpdfstring{\pa{ਸਬਜ਼ੀ} (sabzi)}{ਸਬਜ਼ੀ (sabzi)}]{One noun for a vegetable or cooked vegetable dish}

\label{lesson:PA-C200-sabzi}



\begin{quote}
Retrieve the familiar first and last words of the identity frame.
\end{quote}

\begin{grammarlens}[title={Grammar Lens: direct owner}]
\begin{quote}
\textbf{\pa{ਸਬਜ਼ੀ}} — \emph{sabzī} — vegetable; cooked vegetable dish
\end{quote}

Hear and say the whole word before reading it. Keep it as one whole noun rather than inventing smaller Punjabi pieces. This lesson owns the word alone; the complete identity frame comes only after a separate retrieval.
\end{grammarlens}

\begin{cousinweb}
Punjabi University's Punjabipedia records the present-day Punjabi vegetable category. It supports the bounded everyday vegetable meaning taught here. Source: \href{https://punjabipedia.org/topic.aspx?txt=\%E0\%A8\%B8\%E0\%A8\%AC\%E0\%A9\%9B\%E0\%A9\%80\%E0\%A8\%86\%E0\%A8\%82}{Punjabi University Punjabipedia}.

The exact origin of modern written \textbf{\pa{ਸਬਜ਼ੀ}} is unresolved in the sources used here. Its origin remains unknown here. This lesson therefore makes no etymological claim. Keep the present-day Punjabi noun whole rather than turning spelling resemblance into a derivation. A Government of Punjab canteen rate list provides a current food-service use of the singular noun. Source: \href{https://punjab.gov.in/wp-content/uploads/2022/04/Hospitality-Food-and-Beverages-Rate-List-M.L.A-Hostel-Canteen-Punjab-Chd.pdf}{Government of Punjab}.
\end{cousinweb}

\subsection*{Your turn}
\begin{itemize}
\raggedright
  \item \emph{Hear:} \emph{sabzī}; YOU SAY: \textbf{\pa{ਸਬਜ਼ੀ}}
  \item \emph{Read it:} \textbf{\pa{ਸਬਜ਼ੀ}}
  \item \emph{Write it:} copy \textbf{\pa{ਸਬਜ਼ੀ}} once, checking the whole word
\end{itemize}

\begin{tcolorbox}[breakable,colback=teal!4,colframe=teal!35!black,title={Before you move on}]
Produce the noun for “vegetable” or “cooked vegetable dish” without looking.
\end{tcolorbox}
\section[\texorpdfstring{\pa{ਸਬਜ਼ੀ} (sabzi)}{ਸਬਜ਼ੀ (sabzi)}]{Retrieve the new word on its own}

\label{lesson:PA-C200-sabzi-review}



\begin{quote}
\emph{Hear:} \emph{sabzī}; YOU SAY: \textbf{\pa{ਸਬਜ਼ੀ}}
\end{quote}

\subsection*{Your turn}
\begin{itemize}
\raggedright
  \item \emph{Hear:} \emph{vegetable or cooked vegetable dish}; YOU SAY: \textbf{\pa{ਸਬਜ਼ੀ}}
  \item \emph{Read it:} \textbf{\pa{ਸਬਜ਼ੀ}}
  \item \emph{Write it:} cover and write \textbf{\pa{ਸਬਜ਼ੀ}}
\end{itemize}

\begin{grammarlens}[title={Grammar Lens}]
Keep one noun and the bounded vegetable-or-cooked-dish meaning here. Wider forms, compounds, and specialised classifications stay outside this lesson.
\end{grammarlens}

\begin{tcolorbox}[breakable,colback=teal!4,colframe=teal!35!black,title={Before you move on}]
Say and then write the word once from memory.
\end{tcolorbox}
\section[\texorpdfstring{\pa{ਇਹ} \pa{ਸਬਜ਼ੀ} \pa{ਹੈ} (eh sabzi hai)}{ਇਹ ਸਬਜ਼ੀ ਹੈ (eh sabzi hai)}]{Put the new noun into one identity frame}

\label{lesson:PA-C200-vegetable-label}



\begin{quote}
Retrieve \textbf{\pa{ਇਹ}}, then the new noun, then \textbf{\pa{ਹੈ}}.
\end{quote}

\begin{grammarlens}[title={Grammar Lens}]
\begin{quote}
\textbf{\pa{ਇਹ} \pa{ਸਬਜ਼ੀ} \pa{ਹੈ}} — \emph{eh sabzī hai} — this is a vegetable / cooked vegetable dish
\end{quote}

Keep the new whole noun between familiar \textbf{\pa{ਇਹ}} and \textbf{\pa{ਹੈ}}. Familiar \textbf{\pa{ਹੈ}} stays last in the already-owned identity frame. This chapter teaches only this one label; it does not add a form rule or a wider construction.
\end{grammarlens}

\subsection*{Your turn}
\begin{itemize}
\raggedright
  \item \emph{Hear:} \emph{eh sabzī hai}; YOU SAY: \textbf{\pa{ਇਹ} \pa{ਸਬਜ਼ੀ} \pa{ਹੈ}}
  \item \emph{Read it:} \textbf{\pa{ਇਹ} \pa{ਸਬਜ਼ੀ} \pa{ਹੈ}}
  \item \emph{Write it:} cover and write \textbf{\pa{ਇਹ} \pa{ਸਬਜ਼ੀ} \pa{ਹੈ}}
\end{itemize}

\begin{tcolorbox}[breakable,colback=teal!4,colframe=teal!35!black,title={Before you move on}]
Produce the complete identity frame.
\end{tcolorbox}
\section[\texorpdfstring{\pa{ਇਹ} \pa{ਸਬਜ਼ੀ} \pa{ਹੈ} (eh sabzi hai)}{ਇਹ ਸਬਜ਼ੀ ਹੈ (eh sabzi hai)}]{A bounded vegetable-noun payoff}

\label{lesson:PA-C200-payoff}



\begin{quote}
Retrieve the new noun alone before rebuilding its identity frame.
\end{quote}

\subsection*{Your turn: four skills}
\begin{itemize}
\raggedright
  \item \emph{Hear:} \emph{eh sabzī hai}; YOU SAY: \textbf{\pa{ਇਹ} \pa{ਸਬਜ਼ੀ} \pa{ਹੈ}}
  \item \emph{Read it:} \textbf{\pa{ਇਹ} \pa{ਸਬਜ਼ੀ} \pa{ਹੈ}}
  \item \emph{Write it:} cover the line and write \textbf{\pa{ਇਹ} \pa{ਸਬਜ਼ੀ} \pa{ਹੈ}}
\end{itemize}

\begin{tcolorbox}[breakable,colback=teal!4,colframe=teal!35!black,title={Before you move on}]
Say the new noun, then produce the complete identity frame.

This is the twenty-third prerequisite tranche exposed by Punjabi A2 mock drafting. It does not unblock Punjabi A2 mock 1 by itself. It is not an A2 mock, not A2 readiness or attainment, and not external or human validation.
\end{tcolorbox}
